\documentclass[a4paper,12pt]{article}

\usepackage[utf8]{inputenc}
\usepackage[T1]{fontenc}
\usepackage[francais]{babel}
\usepackage{mathptmx}
\usepackage[scaled]{helvet}
%\usepackage{courier}

\usepackage[margin=2.2cm,nohead,foot=1.8cm]{geometry}
\usepackage{parskip}
\usepackage{comment}
\usepackage{graphicx}
\usepackage{fancyvrb}

\begin{document}

\begin{flushleft}
\begin{tabular}[t]{l}
\Large\textsf{\textbf{Rapport sur le manuscrit d'Habilitation à Diriger les Recherches}}\\[4mm]
\Large\textsf{par Monsieur \textbf{Heithem Abbes}}\\[4mm]
\large\textsf{\textbf{pour l'obtention de l'Habilitation à Diriger les Recherches}}\\
\large\textsf{\textbf{de l'Université Paris 13}}\\
\end{tabular}%

\vspace{4mm}

\noindent%
\begin{tabular}[t]{l}
Rapporteur: Stéphane \textsc{Genaud},\\
Professeur des Universités,\\
Université de Strasbourg.\\[5mm]
Titre du document évalué :\\\\[3mm]
\end{tabular}%

\begin{quote}
\Large \textit{\textsf{\textbf{
	  Contributions à la Gestion Décentralisée des Ressources dans les Environnements de Grilles et Cloud
}}}
\end{quote}

\end{flushleft}

\vspace{3mm} Le document  décrit les travaux de recherche  réalisés par Monsieur
Heithem Abbes au laboratoire de recherche LaTICE, à Tunis, en collaboration avec
le LIPN de  l'Université de Paris 13,  depuis sa thèse de doctorat  en 2009. Ces
activités ont pour  objectif d'améliorer la conception et la  mise en oeuvre des
intergiciels de gestion de ressources de calcul et de données.

\section*{Analyse du document}
\label{sec-1}
Le  document  est organisé  en  cinq  chapitres.  Le  premier chapitre  est  une
introduction qui donne le contexte scientifique  dans lequel se situe ce travail
et résume les principales contributions.

Le  chapitre 2  brosse ensuite  l'évolution, dans  les dernières  décennies, des
systèmes distribués  de calcul,  sous le prisme  des concepts  et problématiques
relatifs au travail de recherche présenté.


Le chapitre 3  décrit une première série de contributions  concernant la gestion
des  données dans  des environnements  décentralisés.  Il  introduit d'abord  la
problématique  de la  gestion  des données  volumineuses  pour des  applications
scientifiques de  type workflow.  Le chapitre  décrit avec du recul,  le premier
travail de l'auteur dans ce domaine, l'intergiciel \emph{BonjourGrid} (2008). Il
fait le constat que la gestion  des données de BonjourGrid devait être améliorée
pour que les applications manipulant  des données volumineuses puissent passer à
l'échelle.  Cette fonctionnalité initialement traitée comme une tâche auxiliaire
du  coordinateur,  fait  l'objet  d'une  refonte  pour  devenir  une  composante
\emph{gestionnaire  de données}  à part  entière de  l'intergiciel.  De  manière
pertinente, l'auteur a choisi de s'appuyer  sur des projets existants (BitDew et
Stork) pour réaliser  cette composante. Disposer d'un gestionnaire  de données a
permis de  concevoir des solutions  de cache  et de pré-chargement  des données,
indépendantes de l'intergiciel (l'auteur l'a intégré dans BonjourGrid mais aussi
dans Condor). Le calcul du pré-chargement de données, étant donné un workflow de
tâches (exprimant les  dépendances) est un problème tout à  fait général, que le
document aurait pu mettre davantage en  exergue. Ces travaux sur le gestionnaire
de  données  ont  fait  l'objet d'évaluations  expérimentales  conséquentes  (de
l'ordre  de 200  machines) sur  une  plate-forme réelle,  Grid'5000.  C'est  une
démarche  qu'il faut  saluer  car elle  demande  beaucoup de  temps  de mise  au
point. Elle me semble cependant nécessaire pour avoir une validation crédible de
la proposition.

Toujours en matière de gestion décentralisée des données, ce chapitre expose des
travaux sur la  sécurisation des données dans le cadre  d'un framework MapReduce
par un  mécanisme de  dissémination des  données. La  motivation aurait  pu être
donnée au début de l'exposé (et non au sein du descriptif du fonctionnement) car
cette problématique  est d'importance quand  on fait  appel à des  ressources de
calcul et stockage externes.  En effet, il est ainsi possible  de ne confier aux
entités externes de stockage qu'une  partie des données, rendant impossible leur
compréhension  sans les  parties complémentaires  hébergées sur  des entités  de
confiance.  Cette problématique  est  complètement en  phase avec  l'utilisation
grandissante des  clouds publics  pour lesquels  la confidentialité  des données
hébergées  est   une  question  cruciale.    Pour  ce  travail   également,  des
expérimentations conséquentes sur Grid'5000 ont permis  de mesurer le coût de la
surcouche de dissémination.


%\section{Ch 4 Modèles décentralisés de gestion des points de reprises de conteneurs}
%\label{sec-2}

Le  chapitre  4  traite  de  la   tolérance  aux  pannes  par  un  mécanisme  de
\emph{check-point  and  restart}  pour  des  applications  s'exécutant  dans  un
environnement de  \emph{containers}. Ce thème de  recherche est particulièrement
d'actualité avec le développement des  datacenters, et de nombreux chercheurs de
la recherche académique ou industrielle  y travaillent. Ce chapitre comporte une
description des  travaux relatifs tout  à fait conséquente et  précise, montrant
bien  comment se  situe la  recherche de  l'auteur. L'originalité  de l'approche
présentée  ici  réside  dans  le  stockage des  points  de  reprise  de  manière
décentralisée (alors que ce stockage  est traditionnellement fait sur un support
centralisé  réputé fiable)  ainsi que  l'ajout d'un  mécanisme de  \emph{garbage
  collector},  qui élimine  au  fur et  à mesure  les  checkpoints périmés.   La
distribution  des check-points  dans le  système est  opérée avec  une table  de
hashage  distribuée.  Cette  table  de  hashage repose  sur  un logiciel  ouvert
éprouvé (Cassandra), modifié  pour s'adapter à la gestion des  snapshots et à la
détection  de pannes  des  hôtes. L'approche  est  difficile techniquement  mais
ambitieuse car  la sauvegarde  de containers offre  une solution  générique pour
tout type  d'application.  Pour ce  travail encore, des  expérimentations menées
grâce à Grid'5000 renforcent considérablement la validité de la proposition.



Le dernier  chapitre du manuscrit  est consacré aux conlusions  et perspectives.
Les perspectives font  état des pistes à court terme,  qui visent l'amélioration
du mécanisme  de tolérance aux  pannes par \emph{check-point and  restart}.  Une
part conséquente du chapitre aborde également  des perspectives de moyen et long
termes. Les idées présentées font apparaître,  en particulier pour ce qui est de
l'utilisation  de navigateurs  comme  noeud  de calcul  et  stockage, une  vraie
conviction personnelle et enthousiasmante de l'auteur.



\section*{Conclusion}


Le manuscrit  retrace les travaux de  recherche menés par Monsieur  Heithem Abbes
depuis sa thèse  ainsi que pendant les  dix années qui ont suivies.   Lors de sa
thèse de doctorat, effectuée en  co-tutelle entre l'Ecole Nationale des Sciences
de l'Informatique à Tunis  et le LIPN à l'Université Paris  13, Monsieur Abbes a
pleinement contribué à  construire une collaboration fructueuse  et pérenne avec
les chercheurs  français.  Il est  remarquable que  Monsieur Abbes ait  réussi à
animer  cette  collaboration  durant  ces   années  en  dépit  de  l'éloignement
géographique  et   des  charges   lourdes  d'une  fonction   d'enseignant.   Ses
publications pendant cette  période montrent que le candidat a  su faire évoluer
ses  recherches  avec  l'état de  l'art  du  domaine,  tout  en gardant  le  fil
conducteur  de la  gestion décentralisée  des  traitements et  données dans  les
systèmes distribués.   Les publications montrent également  que dans l'évolution
de ses recherches, il a su faire  naitre de nouvelles collaborations tout à fait
naturelles avec d'autres  chercheurs français.  Il a  également co-encadré trois
doctorants tunisiens.

Au final,  Monsieur Heithem Abbes fait  dans son manuscrit une  bonne synthèse de
son  activité de  recherche  depuis  son doctorat.  Il  fait  apparaitre un  fil
conducteur  autour duquel  il  résume  avec recul  les  points importants  d'une
recherche  ayant  évolué   avec  les  problématiques  les   plus  actuelles  qui
intéressent la communauté. Ce parcours de recherche a donné lieu à de nombreuses
publications, dont 9 articles dans des revues internationales et une quinzaine de
publications dans des conférences internationales de haut niveau.

Pour toutes les raisons évoquées ci-dessus, je suis favorable à la soutenance de
son habilitation à diriger les recherches.


%\section{Questions}
%\label{sec-4}
%\begin{itemize}
%\item Q chap4 : sauvegarde differentielle des checkpoints ?
%\end{itemize}


\vspace{.8cm}
\begin{flushright}
Stéphane \textsc{Genaud},\\
{\small fait à Strasbourg, le 20 novembre 2019}.\\
\includegraphics[width=.20\textwidth]{../signgenaud.jpg}
\end{flushright}
\end{document}
